\clearpage
\SourceLeaf{445}
\fancyhead[L]{\small STÖRUNGEN DER PALLAS. NACHLASS.}
\pdfbookmark[1]{Nachlass. [II.] Exposition d'une nouvelle méthode}{pallas-II-exposition}
\begin{center}
{\Large NACHLASS.\par}
\bigskip
{\large [II.] EXPOSITION\par}
{\large D’UNE NOUVELLE MÉTHODE\par}
DE CALCULER LES PERTURBATIONS PLANÉTAIRES\par
AVEC L’APPLICATION AU CALCUL NUMÉRIQUE\par
DES PERTURBATIONS DU MOUVEMENT DE PALLAS.
\end{center}
\begin{flushright}
\begin{minipage}{.94\textwidth}
\small\raggedright
. . . . . their motions, periods and their laws\\
Give me to scan.\quad \textsc{Thomsons} Autumn 1346.
\par\medskip
See, how associate round their central sun,\\
Their faithful rings the circling planets run\\
Exactly tracing their appointed sphere.\\
\hspace*{1em}Boyle on deity: Elegant Extracts p. 126.
\end{minipage}
\end{flushright}
\noindent\hspace*{\parindent}Le calcul des perturbations planétaires est susceptible de deux formes différentes. Sous l’une de ces formes, on rapporte le mouvement d’une planète à des élémens constans, ou affectés seulement de variations séculaires, en corrigeant le rayon vecteur, la longitude dans l’orbite et la latitude calculées dans cette hypothèse, par des équations périodiques; sous l’autre forme le mouvement est rapporté a une ellipse variable, laquelle dans chaque instant, sans autre correction, représente exactement le mouvement. On connait assez le succès, avec lequel on s’est servi jusqu’ici exclusivement de la première forme pour les perturbations des anciennes planètes. Cependant il est reconnu, que ce succès dépend surtout de plusieurs circonstances favorables, qui ne sont plus les mêmes pour quelques unes des nouvelles planètes, l’orbite de Pallas étant inclinée sur celle de Jupiter de 34 degrés avec une \SourceLeaf{446}excentricité, qui est presque égale à \(\frac{1}{4}\), et l’orbite de Junon, avec une inclinaison très-considérable, étant encore plus excentrique. Dans ces cas les anciennes méthodes ont paru impraticables.
\par Les nouvelles méthodes, qui font l’objet de ce mémoire, présentent originairement les perturbations sous la seconde forme. Il serait en effet facile, de transformer les perturbations des élémens, une fois trouvées, en perturbations des trois coordonnées, si on le jugeait à propos: cependant nous sommes d’avis, qu’on peut bien se dispenser de ce travail, et s’en tenir aux perturbations des élémens. Quoique ces derniers soient au nombre de six, et que les perturbations des trois coordonnées ne semblent exiger dans les applications numériques que la moitié du travail, ce désavantage n’est pas réel toutes les fois qu’il s’agit de calculer plusieurs lieux de la planète séparés par un modique intervalle de tems, puisque dans ce cas, qui est le plus fréquent dans la pratique, on peut fort bien s’en tenir à un seul système d’élémens, tandisque l’autre méthode exigera du moins le calcul des inégalités des coordonnées pour deux époques différentes. La commodité d’une table auxiliaire pour l’équation du centre et pour le rayon vecteur dans une ellipse constante, à laquelle il faut renoncer en rapportant toutes les inégalités aux élémens mêmes, est aussi fort peu de chose, lorsqu’on opère sur une planète sujette à autant d’inégalités comme Pallas, et peut même être compensée par d’autres simplifications de calcul qu’il serait superflu de rappeler ici.
\par\bigskip\centerline{\rule{3cm}{.4pt}}\par
